\begin{theorem}[Quantitative obstruction-chain candidate]
Across the current TwistedMerge evidence table, residual diagnostics enter the
merge pipeline through distinct decision channels rather than through one
universal scalar obstruction.  In the planted benchmark, cycle score predicts
pairwise merge degradation with Spearman correlation 0.796
over 80 rows.  Validation deltas for selector rows have
Spearman correlation 0.552 with test selector gain over
948 rows, but this is selection evidence rather than a residual
certificate.  Projection residual gates have accepted-row false lift rate
0.000.  Period-index rows match selected
lift decisions at rate 1.000.  Finally, the block-gauge
phase diagram contains 12240 rows where projected cycle score
alone would be misleading without the connection-residual gate.
\end{theorem}

\begin{remark}
The theorem candidate is empirical and artifact-scoped.  It supports the claim
that different residual scores predict different downstream decisions.  It does
not support the stronger claim that every merging failure is Brauer/projective.
\end{remark}
